\documentclass[10pt,a4paper]{amsart}
\usepackage{amsmath,amssymb,amsthm,graphicx}
% For commutative diagrams we use Xy-pic
\usepackage[all]{xy}

% We use 2cell for 2-commutative diagrams.
\xyoption{2cell}
% Theorem environments.
%
\theoremstyle{plain}
\newtheorem{theorem}[subsection]{Theorem}
\newtheorem{proposition}[subsection]{Proposition}
\newtheorem{lemma}[subsection]{Lemma}

\theoremstyle{definition}
\newtheorem{definition}[subsection]{Definition}
\newtheorem{example}[subsection]{Example}
\newtheorem{exercise}[subsection]{Exercise}
\newtheorem{situation}[subsection]{Situation}

\theoremstyle{remark}
\newtheorem{remark}[subsection]{Remark}
\newtheorem{remarks}[subsection]{Remarks}

\numberwithin{equation}{subsection}

% Macros
%
\def\lim{\mathop{\mathrm{lim}}\nolimits}
\def\colim{\mathop{\mathrm{colim}}\nolimits}
\def\Spec{\mathop{\mathrm{Spec}}}
\def\Hom{\mathop{\mathrm{Hom}}\nolimits}
\def\Ext{\mathop{\mathrm{Ext}}\nolimits}
\def\SheafHom{\mathop{\mathcal{H}\!\mathit{om}}\nolimits}
\def\SheafExt{\mathop{\mathcal{E}\!\mathit{xt}}\nolimits}
\def\Sch{\mathit{Sch}}
\def\Mor{\mathop{\mathrm{Mor}}\nolimits}
\def\Ob{\mathop{\mathrm{Ob}}\nolimits}
\def\Sh{\mathop{\mathit{Sh}}\nolimits}
\def\NL{\mathop{N\!L}\nolimits}
\def\CH{\mathop{\mathrm{CH}}\nolimits}
\def\proetale{{pro\text{-}\acute{e}tale}}
\def\etale{{\acute{e}tale}}
\def\QCoh{\mathit{QCoh}}
\def\Ker{\mathop{\mathrm{Ker}}}
\def\Im{\mathop{\mathrm{Im}}}
\def\Coker{\mathop{\mathrm{Coker}}}
\def\Coim{\mathop{\mathrm{Coim}}}

% Boxtimes
%
\DeclareMathSymbol{\boxtimes}{\mathbin}{AMSa}{"02}

%
% Macros for moduli stacks/spaces
%
\def\QCohstack{\mathcal{QC}\!\mathit{oh}}
\def\Cohstack{\mathcal{C}\!\mathit{oh}}
\def\Spacesstack{\mathcal{S}\!\mathit{paces}}
\def\Quotfunctor{\mathrm{Quot}}
\def\Hilbfunctor{\mathrm{Hilb}}
\def\Curvesstack{\mathcal{C}\!\mathit{urves}}
\def\Polarizedstack{\mathcal{P}\!\mathit{olarized}}
\def\Complexesstack{\mathcal{C}\!\mathit{omplexes}}
% \Pic is the operator that assigns to X its picard group, usage \Pic(X)
% \Picardstack_{X/B} denotes the Picard stack of X over B
% \Picardfunctor_{X/B} denotes the Picard functor of X over B
\def\Pic{\mathop{\mathrm{Pic}}\nolimits}
\def\Picardstack{\mathcal{P}\!\mathit{ic}}
\def\Picardfunctor{\mathrm{Pic}}
\def\Deformationcategory{\mathcal{D}\!\mathit{ef}}
\usepackage{cleveref}
\crefname{originaldefinition}{original definition}{original definitions}
\Crefname{originaldefinition}{Original definition}{Original definitions}
\makeatletter
\newlabel{situation-setup}{{02QL}{source}{}{sourcepointer-1}{}}
\newlabel{situation-setup@cref}{{[originaldefinition][0][]02QL}{[1][0][]source}{}{}{}}
\makeatother
\crefname{originallemma}{original lemma}{original lemmas}
\Crefname{originallemma}{Original lemma}{Original lemmas}
\makeatletter
\newlabel{lemma-locally-equidimensional}{{0FE3}{source}{}{sourcepointer-2}{}}
\newlabel{lemma-locally-equidimensional@cref}{{[originallemma][0][]0FE3}{[1][0][]source}{}{}{}}
\makeatother
\crefname{originalremark}{original remark}{original remarks}
\Crefname{originalremark}{Original remark}{Original remarks}
\makeatletter
\newlabel{remark-fundamental-class}{{0FE4}{source}{}{sourcepointer-3}{}}
\newlabel{remark-fundamental-class@cref}{{[originalremark][0][]0FE4}{[1][0][]source}{}{}{}}
\makeatother
\crefname{originalsection}{original section}{original sections}
\Crefname{originalsection}{Original section}{Original sections}
\makeatletter
\newlabel{algebra-section-smooth-overview}{{00TZ}{source}{}{sourcepointer-4}{}}
\newlabel{algebra-section-smooth-overview@cref}{{[originalsection][0][]00TZ}{[1][0][]source}{}{}{}}
\makeatother
\crefname{originallemma}{original lemma}{original lemmas}
\Crefname{originallemma}{Original lemma}{Original lemmas}
\makeatletter
\newlabel{algebra-lemma-regular-ring-CM}{{00NQ}{source}{}{sourcepointer-5}{}}
\newlabel{algebra-lemma-regular-ring-CM@cref}{{[originallemma][0][]00NQ}{[1][0][]source}{}{}{}}
\makeatother
\crefname{originallemma}{original lemma}{original lemmas}
\Crefname{originallemma}{Original lemma}{Original lemmas}
\makeatletter
\newlabel{algebra-lemma-CM-goes-up}{{045J}{source}{}{sourcepointer-6}{}}
\newlabel{algebra-lemma-CM-goes-up@cref}{{[originallemma][0][]045J}{[1][0][]source}{}{}{}}
\makeatother
\crefname{originallemma}{original lemma}{original lemmas}
\Crefname{originallemma}{Original lemma}{Original lemmas}
\makeatletter
\newlabel{lemma-gysin-easy}{{0FEA}{source}{}{sourcepointer-7}{}}
\newlabel{lemma-gysin-easy@cref}{{[originallemma][0][]0FEA}{[1][0][]source}{}{}{}}
\makeatother
\title{The Gysin map of a local complete intersection morphism sends the fundamental class of the target to the fundamental class of the source}
\author{The Stacks Project Authors}
\date{stacks-project:chow.tex:local-snapshot-20261012; GNU FDL 1.2 or later}
\begin{document}
\maketitle
\noindent\textit{Editorial note.} All delivered mathematical text is the author's own, copied from the pinned Stacks Project chapter file chow.tex; no statement or proof step was rewritten, reordered or removed.\par\medskip
\noindent\textit{Editorial note.} Preamble provenance: the pinned chapter file opens with an input of the shared preamble and carries no preamble text of its own, so the project's single Stacks-wide preamble staged at sources/stacks-chow/preamble.tex is excerpted verbatim as two spans, lines 14-18 (Xy-pic setup) and lines 154-217 (the native theorem declarations and mathematical macros). Nothing was added, removed or redefined; the wrapper class and the amsmath/amssymb/amsthm/graphicx base packages are supplied by the assembler.\par\medskip
\noindent\textit{Editorial note.} The counted statement and proof contain no cross reference, no citation and no included file, so the delivered body is self contained.\par\medskip
\noindent\textit{Editorial note.} Canonical location of the counted proposition: Stacks tag 0FF7 (sources/stacks-chow/pinned-original-source.tex line 13643).\par\medskip
\noindent\textit{Allowed standard prerequisites (editorial).}\begin{itemize}
\item The base situation (S, delta) and its bivariant operational Chow groups, together with the fundamental class [Y'] of a scheme. Same chapter, Situation labelled situation-setup (source line 1372) and Remark labelled remark-fundamental-class (source line 8405), both invoked in the statement. Fix the ambient situation and the meaning of the fundamental classes [Y'] and [X'] as sums of integral cycles over the locally constant dimension functions; the reduction of those sums is carried out in the counted span.
\item The equidimensionality conditions and the Gysin statement for a regular closed immersion. Same chapter, Lemma labelled lemma-locally-equidimensional (source line 8277) and Lemma labelled lemma-gysin-easy (source line 12128). Supply the conditions listed as hypothesis (2) of the statement and the final step in the regular closed immersion case at proof line 13717.
\item Regular local rings are Cohen-Macaulay, smoothness gives a flat map with regular fibres, and Cohen-Macaulayness goes up along such a map. Algebra, Section labelled algebra-section-smooth-overview (Stacks tag 00TZ), Lemma labelled algebra-lemma-regular-ring-CM (Stacks tag 00NQ) and Lemma labelled algebra-lemma-CM-goes-up (Stacks tag 045J), invoked at proof lines 13704-13709. Supply the three inputs that let the Cohen-Macaulay hypothesis (4) be transported to the base-changed immersion; the factorization of f and the construction of i' are inside the counted span.
\end{itemize}\medskip
\ifdefined\hypertarget\hypertarget{sourcepointer-1}{}\fi
\noindent\textit{Source reference.} Original definition 02QL: sources/stacks-chow/pinned-original-source.tex, source lines 1372--1372. Same-chapter pointer into the author environment situation at source line 1372; invoked at proof line None of the counted proof. It is retained only as a navigation pointer; every load-bearing step is inside the delivered span.\par\medskip
\ifdefined\hypertarget\hypertarget{sourcepointer-2}{}\fi
\noindent\textit{Source reference.} Original lemma 0FE3: sources/stacks-chow/pinned-original-source.tex, source lines 8277--8277. Same-chapter pointer into the author environment lemma at source line 8277; invoked at proof line None of the counted proof. It is retained only as a navigation pointer; every load-bearing step is inside the delivered span.\par\medskip
\ifdefined\hypertarget\hypertarget{sourcepointer-3}{}\fi
\noindent\textit{Source reference.} Original remark 0FE4: sources/stacks-chow/pinned-original-source.tex, source lines 8405--8405. Same-chapter pointer into the author environment remark at source line 8405; invoked at proof line None of the counted proof. It is retained only as a navigation pointer; every load-bearing step is inside the delivered span.\par\medskip
\ifdefined\hypertarget\hypertarget{sourcepointer-4}{}\fi
\noindent\textit{Source reference.} Original section 00TZ: sources/stacks-chow/tags.txt, source lines 1093--1093. Cross-chapter standard prerequisite located in the project tag list staged in the local snapshot; invoked at proof line 13704 of the counted proof. It is retained only as a navigation pointer; every load-bearing step is inside the delivered span.\par\medskip
\ifdefined\hypertarget\hypertarget{sourcepointer-5}{}\fi
\noindent\textit{Source reference.} Original lemma 00NQ: sources/stacks-chow/tags.txt, source lines 868--868. Cross-chapter standard prerequisite located in the project tag list staged in the local snapshot; invoked at proof line 13706 of the counted proof. It is retained only as a navigation pointer; every load-bearing step is inside the delivered span.\par\medskip
\ifdefined\hypertarget\hypertarget{sourcepointer-6}{}\fi
\noindent\textit{Source reference.} Original lemma 045J: sources/stacks-chow/tags.txt, source lines 5438--5438. Cross-chapter standard prerequisite located in the project tag list staged in the local snapshot; invoked at proof line 13709 of the counted proof. It is retained only as a navigation pointer; every load-bearing step is inside the delivered span.\par\medskip
\ifdefined\hypertarget\hypertarget{sourcepointer-7}{}\fi
\noindent\textit{Source reference.} Original lemma 0FEA: sources/stacks-chow/pinned-original-source.tex, source lines 12128--12128. Same-chapter pointer into the author environment lemma at source line 12128; invoked at proof line 13717 of the counted proof. It is retained only as a navigation pointer; every load-bearing step is inside the delivered span.\par\medskip
\begin{lemma}
\label{lemma-lci-gysin-easy}
Let $(S, \delta)$ be as in Situation \ref{situation-setup}.
Consider a cartesian diagram
$$
\xymatrix{
X' \ar[d]_{f'} \ar[r] &
X \ar[d]^f \\
Y' \ar[r] &
Y
}
$$
of schemes locally of finite type over $S$. Assume
\begin{enumerate}
\item $f$ is a local complete intersection morphism and
the Gysin map exists for $f$,
\item $X$, $X'$, $Y$, $Y'$ satisfy the equivalent conditions of
Lemma \ref{lemma-locally-equidimensional},
\item for $x' \in X'$ with images $x$, $y'$, and $y$
in $X$, $Y'$, and $Y$ we have $n_{x'} - n_{y'} = n_x - n_y$
where $n_{x'}$, $n_x$, $n_{y'}$, and $n_y$ are as in the lemma, and
\item for every generic point $\xi \in X'$ the local ring
$\mathcal{O}_{Y', f'(\xi)}$ is Cohen-Macaulay.
\end{enumerate}
Then $f^![Y'] = [X']$ where $[Y']$ and $[X']$ are as in
Remark \ref{remark-fundamental-class}.
\end{lemma}

\begin{proof}
Recall that $n_{x'}$ is the common value of $\delta(\xi)$
where $\xi$ is the generic point of an irreducible component
passing through $x'$. Moreover, the functions
$x' \mapsto n_{x'}$, $x \mapsto n_x$, $y' \mapsto n_{y'}$, and
$y \mapsto n_y$ are locally constant. Let $X'_n$, $X_n$, $Y'_n$,
and $Y_n$ be the open and closed subscheme of $X'$, $X$, $Y'$, and
$Y$ where the function has value $n$. Recall that
$[X'] = \sum [X'_n]_n$ and $[Y'] = \sum [Y'_n]_n$.
Having said this, it is clear that to prove the lemma we
may replace $X'$ by one of its connected components
and $X$, $Y'$, $Y$ by the connected component that
it maps into. Then we know that $X'$, $X$, $Y'$, and
$Y$ are $\delta$-equidimensional in the sense that
each irreducible component has the same $\delta$-dimension.
Say $n'$, $n$, $m'$, and $m$ is this common value
for $X'$, $X$, $Y'$, and $Y$. The last assumption
means that $n' - m' = n - m$.

\medskip\noindent
Choose a factorization $f = g \circ i$ where $i : X \to P$
is an immersion and $g : P \to Y$ is smooth. As $X$ is connected,
we see that the relative dimension of $P \to Y$ at points of $i(X)$
is constant. Hence after replacing $P$ by an open neighbourhood
of $i(X)$, we may assume that $P \to Y$ has constant relative dimension
and $i : X \to P$ is a closed immersion.
Denote $g' : Y' \times_Y P \to Y'$ the base change of $g$ and denote
$i' : X' \to Y' \times_Y P$ the base change of $i$.
It is clear that $g^*[Y] = [P]$ and $(g')^*[Y'] = [Y' \times_Y P]$.
Finally, if $\xi' \in X'$ is a generic point, then
$\mathcal{O}_{Y' \times_Y P, i'(\xi)}$ is Cohen-Macaulay.
Namely, the local ring map
$\mathcal{O}_{Y', f'(\xi)} \to \mathcal{O}_{Y' \times_Y P, i'(\xi)}$
is flat with regular fibre
(see Algebra, Section \ref{algebra-section-smooth-overview}),
a regular local ring is Cohen-Macaulay
(Algebra, Lemma \ref{algebra-lemma-regular-ring-CM}),
$\mathcal{O}_{Y', f'(\xi)}$ is Cohen-Macaulay by assumption
(4) and we get what we want from
Algebra, Lemma \ref{algebra-lemma-CM-goes-up}.
Thus we reduce to the case discussed in the next paragraph.

\medskip\noindent
Assume $f$ is a regular closed immersion and $X'$, $X$, $Y'$, and
$Y$ are $\delta$-equidimensional of $\delta$-dimensions
$n'$, $n$, $m'$, and $m$ and $m' - n' = m - n$.
In this case we obtain the result immediately from
Lemma \ref{lemma-gysin-easy}.
\end{proof}

\end{document}
